% ============================================================
%  ch01  第一讲：仿射 Grassmann 簇及其初步性质
%  原文 Lecture I（原文行 727–902）
% ============================================================

\section{第一讲：仿射 Grassmann 簇及其初步性质}\label{Lecture I}

我们先给出将要研究的对象。

\subsection{$\GL_n$ 的仿射 Grassmann 簇}\label{aff GLn}

先要说明什么叫「一族格」。所谓 $k((t))^n$ 中的一族格，是指 $k((t))^n$ 的一族有限生成投射 $R[[t]]$-子模，张量回 $R((t))$ 之后恰好给出整个 $R((t))^n$。

\begin{dfn}
设 $R$ 是 $k$-代数。所谓 $k((t))^n$ 中的一个 \textbf{$R$-族格}，是指 $R((t))^n$ 的一个有限生成投射 $R[[t]]$-子模 $\La$，使得
\[
\La\otimes_{R[[t]]}R((t))=R((t))^n .
\]
\end{dfn}

\begin{dfn}
（$\GL_n$ 的）\textbf{仿射 Grassmann 簇} $\Gr_{\GL_n}$ 是这样的预层：它给每个 $k$-代数 $R$ 配上 $k((t))^n$ 中全体 $R$-族格构成的集合。
\end{dfn}

为简便起见，本小节把 $\Gr_{\GL_n}$ 简记作 $\Gr$。整个理论都建立在下面这个定理之上。

\begin{thm}\label{rep}
仿射 Grassmann 簇 $\Gr$ 可由一个归纳射影概形表示。
\end{thm}

回忆一下：这尤其意味着 $\Gr$ 可以写成一列有限维射影概形的递增并（见 \S\ref{Intro:Indsch}）。

定理的证明略有技术性，但背后的想法却非常简单，所以我们先做一个非正式的说明（只在 $k$-点层面）。取标准格 $\La_0=k[[t]]^n$。以 $\Gr^{(N)}\subset\Gr$ 记分类那些夹在 $t^N\La_0\subset\La\subset t^{-N}\La_0$ 之间的格 $\La$ 的子空间，则 $\Gr$ 是这些 $\Gr^{(N)}$ 的递增并。于是只需证明 $\Gr^{(N)}$ 可由射影概形表示。

对格 $\La\in\Gr^{(N)}$，商 $\La/t^N\La_0$ 可以看作 $k^{2nN}\cong t^{-N}\La_0/t^N\La_0$ 的子空间，并且它在 $t$ 的作用下稳定。这样一来，$\Gr^{(N)}$ 就实现为通常 Grassmann 簇 $\Gr(2nN)$ 的闭子概形，而后者分类 $k^{2nN}$ 的有限维子空间；因此 $\Gr^{(N)}$ 是射影概形。下面把这个想法写成严格证明。

\begin{proof}
给定一族 $R$-格 $\La\subset R((t))^n$，总存在足够大的整数 $N$，使得
\begin{equation}\label{subfun}
t^NR[[t]]^n\subset \La\subset t^{-N}R[[t]]^n .
\end{equation}
所以 $\Gr$ 是满足 \eqref{subfun} 的那些格所成子函子 $\Gr^{(N)}$ 的并。关键的观察是下面这个引理。

\begin{lem}\label{projectivity}
通过映射
\[
\La\mapsto Q=t^{-N}R[[t]]^n/\La ,
\]
$\Gr^{(N)}$ 与下述预层 $\Gr^{(N),f}$ 等同：
\[
\Gr^{(N),f}:=\left\{Q\ \middle|\ \text{$Q$ 是 $\dfrac{t^{-N}R[[t]]^n}{t^NR[[t]]^n}$ 的商模，且作为 $R$-模投射}\right\}.
\]
\end{lem}

后面会看到，这个论断对 $k$-点来说相当显然，对一般的 $R$ 却颇为微妙。我们还需要一个引理。

以 $\Gr(m)$ 记通常的 Grassmann 簇，它分类 $k^m$ 的有限维子空间，是各个 $\Gr(r,m)$ 的无交并，后者分类 $r$ 维子空间。

\begin{lem}\label{Embedding}
函子 $\Gr^{(N),f}$ 可由 $\Gr(2nN)$ 的一个闭子概形表示。
\end{lem}

先证引理 \ref{Embedding}。把 $t^{-N}k[[t]]^n/t^Nk[[t]]^n$ 与 $k^{2nN}$ 等同，则乘以 $t$ 给出该向量空间上的幂零算子 $\Phi$，而 $\Gr^{(N),f}$ 就是 $\Gr(2nN)$ 中分类那些在 $\Phi$ 下稳定的子空间的闭子概形。注意 $\Gr^{(N),f}$ 实际上是一些 Springer 纤维的并（\cite[\S~1]{Yu}）。

再证引理 \ref{projectivity}。首先要说明引理中的映射
\begin{equation}\label{I:map}
\Gr^{(N)}\to \Gr^{(N),f}
\end{equation}
确实有定义，即 $Q$ 作为 $R$-模投射。由于 $\Lambda/t\Lambda$ 是 $R$-投射的，
\[
R((t))^n/\Lambda\cong\bigoplus_{k\geq 0}t^{-k-1}\Lambda/t^{-k}\Lambda
\]
也 $R$-投射。于是由 $R$-模的短正合列
\[
0\to Q\to R((t))^n/\Lambda\to R((t))^n/t^{-N}R[[t]]^n\to 0
\]
即得 $Q$ 的投射性。

因此 \eqref{I:map} 把 $\Gr^{(N)}$ 实现为 $\Gr^{(N),f}$ 的子预层。还需说明 \eqref{I:map} 在每个 $R$ 的 $R$-点上满射。换言之，对每个 $k$-代数 $R$ 以及 $\Gr^{(N),f}$ 的一个由商映射 $t^{-N}R[[t]]^n/t^NR[[t]]^n\to Q$ 表示的 $R$-点，要说明 $R[[t]]$-模
\[
\Lambda:=\ker\bigl( t^{-N}R[[t]]^n\to  t^{-N}R[[t]]^n/t^NR[[t]]^n\to Q\bigr)
\]
是有限投射的。此时 $\Lambda\otimes R((t))\simeq R((t))^n$ 自动成立。

由引理 \ref{Embedding}，$\Gr^{(N),f}$ 局部有限型，故（见 \cite[Tag01ZC]{St}）
\[
\Gr^{(N),f}(\varinjlim R_i)=\varinjlim \Gr^{(N),f}(R_i).
\]
把 $R$ 写成有限生成子 $k$-代数的并，便可假设 $R$ 在 $k$ 上有限生成。

此时由于 $R$ 是 Noether 环，自然环同态 $R[t]\to R[[t]]$ 平坦（见 \cite[Tag00MB]{St}），并诱导同构 $t^{-N}R[t]^n/t^NR[t]^n\cong t^{-N}R[[t]]^n/t^NR[[t]]^n$。于是可定义 $R[t]$-模
\[
\Lambda_f:= \ker\bigl( t^{-N}R[t]^n\to  t^{-N}R[t]^n/t^NR[t]^n\to Q\bigr).
\]
由 $R[t]\to R[[t]]$ 的平坦性，$\Lambda=\Lambda_f\otimes_{R[t]}R[[t]]$。因此只需证明 $\Lambda_f$ 是有限投射 $R[t]$-模。显然 $\Lambda_f$ 有限生成，故只需说明它在 $R[t]$ 上平坦。因 $\La_f$ 已 $R$-平坦，由 \cite[Tag00MH]{St}，只需说明对每个 $x\in \Spec R$（剩余域记为 $K$），基变换 $\Lambda_f\otimes_{R[t]}K[t]$ 在 $K[t]$ 上平坦。由于 $Q$ 是 $R$-平坦的，
\[
\Lambda_f\otimes_{R[t]}K[t]= \ker \bigl( t^{-N}K[t]^n\to Q\otimes_R K\bigr),
\]
它无 $t$-挠，因而在 $K[t]$ 上平坦。引理 \ref{projectivity} 证毕。
\end{proof}

最后，不难看出含入 $\Gr^{(N)}\to \Gr^{(N+1)}$ 是闭嵌入。定理得证。

\begin{rmk}把 $R[[t]]$ 换成 $R[t]$ 这一想法，是 Beauville--Laszlo 用一条整体曲线实现 $\Gr$ 的特例。见下文定理 \ref{glob aff}。
\end{rmk}

\subsection{一般群的仿射 Grassmann 簇}

设 $G$ 是光滑仿射 $k$-群。回忆主丛的概念（见 \S\ref{Intro:torsor}）以及圆盘的 $R$-族（见 \S\ref{Intro:disc}）。$G$ 的\textbf{仿射 Grassmann 簇} $\Gr_{G}$ 定义为
\begin{equation}\label{loc aff}
\Gr_{G}(R)=\left\{(\mE,\beta)\ \left|\ \begin{array}{l}
\mE \text{ 是 } D_R \text{ 上的 } G\text{-主丛},\\
\beta:\mE|_{D^*_R}\cong  \mE^0|_{D^*_R} \text{ 是一个平凡化}
\end{array}\right.\right\}.
\end{equation}

注意这个定义对非常值群概形也说得通。确切地说，若 $\underline G$ 是 $k[[t]]$ 上光滑的仿射群概形，则可完全仿照 \eqref{loc aff} 定义 $\underline G$ 的仿射 Grassmann 簇 $\Gr_{\underline G}$，只须把 $G$ 换成 $\underline G$。

\begin{thm}\label{I:rep aff}
预层 $\Gr_{\underline G}$ 可由 $k$ 上一个局部有限型的归纳概形表示。若 $\underline G$ 约化，则 $\Gr_{\underline G}$ 是归纳射影的。
\end{thm}

\begin{rmk}
(i) 上述定义与定理中当然可以把 $k[[t]]$ 换成任一局部域的整数环。

(ii) 容易看出，对加法群 $G=\bG_a$，$\Gr_{\bG_a}$ 可由仿射空间的归纳极限表示，因此不是归纳射影的。
\end{rmk}

\begin{rmk}在本讲义中，凡是不需要额外代价就能对一般群概形叙述的结果，我们都按一般情形叙述。不熟悉群概形的读者尽可把 $\underline G$ 换成 $k$-群 $G$。
\end{rmk}

定理的证明基于以下两个结果。

\begin{prop}\label{lin rep}
存在忠实的 $k[[t]]$-线性表示 $\rho: \underline G\to \GL_n$，使得 $\GL_n/\underline G$ 拟仿射。此外，若 $\underline G$ 约化，则可选取 $\rho$ 使 $\GL_n/\underline G$ 仿射。
\end{prop}

\begin{proof}第一句见 \cite[Proposition 1.3]{PR}，第二句见 \cite[Corollary 9.7.7]{Ap}。
\end{proof}

\begin{prop}\label{gen G}
设 $\rho:\underline G\to \GL_n$ 是线性表示，且 $\GL_n/\underline G$ 拟仿射，则诱导映射 $f_\rho:\Gr_{\underline G}\to \Gr_{\GL_n}$ 是局部闭嵌入。此外，若 $\GL_n/\underline G$ 仿射，则它是闭嵌入。
\end{prop}

这个证明与后文无关，愿意接受该命题的读者可以跳过。

\begin{proof}
设 $\Spec R\to \Gr_{\GL_n}$ 是一映射，由 $(\mE,\beta)$ 表示。要说明基变换映射
\[
\mF:=\Spec R\times_{\Gr_{\GL_n}}\Gr_{\underline G}\to \Spec R
\]
是局部闭嵌入，并且在 $\GL_n/\underline G$ 仿射时是闭嵌入。以 $\pi:\mE\to D_R$ 记投影，并把 $\beta$ 看作 $\pi$ 在 $D^*_R$ 上的截面。以
\[
\bar\pi:[\mE/\underline G]\to D_R
\]
记 fpqc 商。由于 $\GL_n/\underline G$（拟）仿射，由（拟）仿射态射的忠实平坦下降（见 \cite[Tag0244, Tag0246]{St}），存在 $D_R$ 上仿射且有限展示的概形 $W$，使得 $[\mE/\underline G]\subset W$ 是拟紧开子集。截面 $\beta$ 诱导 $\bar\pi$ 在 $D^*_R$ 上的截面，记作 $\bar\beta$。把 $\mE$ 约化为 $\underline G$-主丛 $\mE'$，等价于给出 $\bar\pi$ 的一个截面。于是 $\mF$ 就是 $\Spec R$ 上的预层：它给每个 $R\to R'$ 配上 $\bar\pi$ 在 $D_{R'}$ 上的截面 $\beta'$，并要求 $\beta'|_{D^*_{R'}}=\bar\beta|_{D^*_{R'}}$。我们有

\begin{lem}
设 $p:V\to D_R$ 是有限展示的仿射概形，$s$ 是 $p$ 在 $D_R^*$ 上的截面。则 $\Spec R$ 上那个预层——它给每个 $R\to R'$ 配上 $p$ 在 $D_{R'}$ 上的截面 $s'$ 且要求 $s'|_{D^*_{R'}}=s|_{D^*_{R'}}$——可由 $\Spec R$ 的闭子概形表示。
\end{lem}

\begin{proof}
把 $V\subset \bA^N_{D_R}$ 嵌入某个（很大的）仿射空间，成为闭子概形。用坐标把 $s$ 写成
\[
s=(s_1(t),\ldots,s_N(t)),\quad s_i(t)=\sum_j s_{ij}t^{j}\in R((t)).
\]
所论预层便由闭子概形 $\Spec A\subset \Spec R$ 表示，其中 $A$ 由 $\{s_{ij}=0,\ i=1,\ldots,N,\ j<0\}$ 生成的理想定义。
\end{proof}

把该引理用于 $V=W$、$s=\bar\beta$，得到闭子概形 $\Spec A\subset\Spec R$，以及一个恒同映射 $\beta_A: D_A\to W_{D_A}$，它在 $D_A^*$ 上的限制是 $\bar\beta|_{D^*_A}$。沿 $t=0$ 给出的闭嵌入 $\Spec A\to D_A$ 作基变换，得到截面 $s_0:\Spec A\to W\times_{D_R}\Spec A$。于是 $\mF$ 由 $\Spec A$ 的开子概形 $s_0^{-1}([\mE/\underline G]\times_{D_R}\Spec A)$ 表示。
\end{proof}

\begin{rmk}\label{general base}
本讲义在固定的基域 $k$ 上工作。但对某些应用（例如 \cite{PR2,PZ}），需要在基概形 $S=\Spec A$ 上工作。从 $A[[t]]$ 上仿射光滑群概形 $\underline G$ 出发，同样可用 \eqref{loc aff} 定义仿射 Grassmann 簇 $\Gr_{\underline G}$，它是 $S$ 上的预层。若相应的命题 \ref{lin rep} 在这种一般性下成立，则与上面相同的论证说明 $\Gr_{\underline G}$ 可表示为 $S$ 上有限展示的归纳概形。已知命题 \ref{lin rep} 在以下情形成立：(i) $\underline G$ 是 $S$ 上线性约化群概形的基变换；(ii) $S$ 是 Dedekind 整环（见 \cite[Appendix]{PZ}）。
\end{rmk}

有意思（也很重要）的非常值群概形，来自 Bruhat--Tits 理论中的抛孔群概形（综述见 \cite{Ti}）。我们不一般地解释它们是什么，只给出两个基本例子。

\begin{ex}\label{hyperspacial and Iwahori}
设 $G$ 是 $k$ 上连通约化群，$B\subset G$ 是 Borel 子群。
\begin{enumerate}
\item 群概形 $\underline G:=G\otimes_k \mO$ 是抛孔的，通常称为 $G$ 的\textbf{超特殊}抛孔。
\item 仍设 $\underline G$ 如上。有求值映射 $\on{ev}: \underline G(\mO)\to G(k)$，$g\mapsto (g\bmod t)$。则存在 $\mO$ 上光滑、纤维连通且仿射的群概形 $\underline G'$，连同同态 $\underline G'\to \underline G$，它诱导同构 $\underline G'\otimes_{\mO}F\cong \underline G\otimes_\mO F=G\otimes_k F$，并把 $\underline G'(\mO)$ 与 $\on{ev}^{-1}(B(k))$ 等同。通常称 $\underline G'$ 为 $G$ 的 \textbf{Iwahori 群概形}。
\end{enumerate}
更多讨论见例 \ref{hyperspacial and Iwahori2}。
\end{ex}

\begin{rmk}\label{terminology}
本讲义把 $\Gr_{\underline G}$ 称为 $\underline G$ 的仿射 Grassmann 簇。这与文献中的标准术语\emph{并不一致}。例如当 $\underline G$ 是如上 $G$ 的 Iwahori 群概形时，文献中通常把 $\Gr_{\underline G}$ 称为 $G$ 的\textbf{仿射旗簇}，并记作 $\Fl_{G}$。
\end{rmk}

下述命题推广了定理 \ref{I:rep aff} 的第二句。

\begin{prop}\label{indproper}
设 $\underline G$ 是 $\mO$ 上纤维连通光滑的仿射群概形，其一般纤维连通约化。则 $\Gr_{\underline G}$ 归纳射影，当且仅当 $\underline G$ 是抛孔群概形。
\end{prop}

这与经典结果类似：连通约化群 $G$ 的齐性空间射影，当且仅当它是 $G$ 的部分旗簇。该结果在本讲义别处不用，故略去证明。

下面引入卷积 Grassmann 簇 $\Gr_{\underline G}\tilde\times\Gr_{\underline G}$，它定义为
\begin{equation}\label{I:conv}
\Gr_{\underline G}\tilde\times\Gr_{\underline G}(R)=\left\{(\mE_1,\mE_2,\beta_1,\beta_2)\ \left|\ \begin{array}{l}
\mE_1,\mE_2 \text{ 是 } D_R \text{ 上的 } \underline G\text{-主丛},\\
\beta_1:\mE_1|_{D_R^*}\cong \mE^0|_{D_R^*},\ \beta_2:\mE_2|_{D_R^*}\cong  \mE_1|_{D_R^*}
\end{array}\right.\right\}.
\end{equation}
有映射
\begin{equation}\label{conv m0}
m:\Gr_{\underline G}\tilde\times\Gr_{\underline G}\to \Gr_{\underline G},\quad (\mE_1,\mE_2,\beta_1,\beta_2)\mapsto (\mE_2, \beta_1\beta_2),
\end{equation}
以及自然投影
\[
\pr_1:\Gr_{\underline G}\tilde\times\Gr_{\underline G}\to \Gr_{\underline G}, \quad (\mE_1,\mE_2,\beta_1,\beta_2)\mapsto (\mE_1,\beta_1),
\]
二者一起诱导同构
\begin{equation}\label{I: isom}
(\pr_1,m):\Gr_{\underline G}\tilde\times\Gr_{\underline G}\cong \Gr_{\underline G}\times \Gr_{\underline G}.
\end{equation}
特别地，卷积 Grassmann 簇可表示。映射 $m$ 通常称为\textbf{卷积映射}。

类似地有 $n$ 重卷积 Grassmann 簇 $\Gr_{\underline G}\tilde\times\cdots\tilde\times\Gr_{\underline G}$，它分类 $(\mE_i,\beta_i)$，$i=1,\ldots,n$，其中 $\mE_i$ 是 $D_R$ 上的 $\underline G$-主丛，$\beta_i: \mE_{i}|_{D_R^*}\simeq \mE_{i-1}|_{D_R^*}$ 是同构（这里 $\mE_0=\mE^0$ 平凡）。对每个 $j=1,\ldots,n-1$，记下 $(\mE_i,\beta_i)$，$i=0,\ldots,j$，便得到到 $j$ 重卷积 Grassmann 簇的自然投影；另有 $n$ 重卷积映射
\begin{equation}\label{conv m}
m:\Gr_{\underline G}\tilde\times\cdots\tilde\times\Gr_{\underline G}\to \Gr_{\underline G},
\end{equation}
它把 $(\mE_i,\beta_i)$ 送到 $(\mE_n, \beta_1\cdots\beta_n)$。
